\chapter{Market Data for Research}\label{ch:rs:market-data-for-research}

Two researchers measure the volatility of the same stock on the same day from one-second returns. One uses trade prices,
the other midquotes, and the first variance is 3.6 times the second. Sampled every five minutes, the two agree within 2\%.
Neither researcher made an error; each chose a representation of the market, and every representation keeps some of what
happened and destroys the rest. This chapter is about those choices: what a message feed records, the clocks by which a
stream of trades is cut into \omterm{def:rs:market-data-for-research:bar}{bars}, what trade prices, midquotes and \omterm{def:rs:market-data-for-research:bar}{bars} each lose, how a continuous history is built from
futures that expire, and how much data all this is. Its data are one simulated trading day from the book's market simulator,
built in this chapter, and nine years of crude oil futures settlements.

\section{What a message feed contains}

One Quant Book 1, chapter 28, described the levels of market data: the best prices (level 1), the depth by price (level 2)
and the book order by order (level 3), with exchange and receive timestamps, sequence numbers and trade conditions. A research
store starts from the richest of these, an order-by-order feed, because every coarser view can be rebuilt from it and no finer
view can be rebuilt from a coarser one.

The book's simulator, \texttt{firm.tape}, writes three messages in the style of Nasdaq's TotalView-ITCH feed (the table
below). An order is added with an identifier, a side, a price and a size; it is later cancelled in whole or part, or executed
against an incoming order. Trades are not separate facts: they are the executions, and the aggressor's side is the opposite of
the resting order's. The simulated day used throughout the chapter has 1\,555\,872 messages and 55\,542 trades, 28 messages a
trade: most messages are quotes placed and withdrawn without ever trading.

\begin{center}
\small
\begin{tabular}{llp{7.6cm}}
\toprule
type & fields & meaning \\
\midrule
\texttt{A} add & time, id, side, price, size & a limit order joins the back of the queue at its price \\
\texttt{X} cancel & time, id, size & part or all of a resting order is withdrawn \\
\texttt{E} execute & time, id, size, aggressor, trade id & a resting order is filled by an incoming order; one row per resting order filled \\
\bottomrule
\end{tabular}
\end{center}

The simulator's mechanism is simple enough to state in a paragraph, and every later chapter that uses it relies on it. An
efficient price, which nobody observes, jumps by a tick or two at random times. Liquidity providers add orders at the ten best
levels of each side and cancel them at a steady rate, faster when the efficient price has moved through their level. Noise
traders send market orders whose arrivals cluster (a Hawkes process, One Quant Book 4, chapter 7) and whose signs persist,
because most are slices of larger parent orders. Informed traders send market orders towards the efficient price when the
mid is far from it. A random activity level, persistent over minutes, and a U-shaped intraday profile scale all of these
rates together, and a news burst in mid-session makes the efficient price move faster for ninety seconds. The simulator
records the truth that real data hide: the efficient price, and which trades were informed.

\begin{dated}{2026-09}{How large is a day of order-by-order data?}\label{dat:rs:market-data-for-research:itch}
Nasdaq publishes sample files of full days of its TotalView-ITCH 5.0 feed on a public server. The compressed (gzip) files for
single days in 2019 and early 2020 are 3.5 to 5.6 billion bytes (30 January 2019: 4.8~GB; 30 January 2020: 5.6~GB); the
files posted between August 2025 and June 2026 are 4.7 to 17.9~GB. That is one exchange's feed, for every stock it trades, for
one day, compressed.
\end{dated}

\section{Sampling clocks}

A strategy that does not trade on every message looks at the market through \omterm{def:rs:market-data-for-research:bar}{bars}: summaries of the trades between two
instants. The instants are chosen by a clock, and the clock need not be the wall clock.

\begin{definition}[Bar, time bar, tick bar, volume bar, dollar bar, VWAP]\label{def:rs:market-data-for-research:bar}
A \emph{bar}\index{bar} summarises a run of consecutive trades by its open, high, low and close prices, its volume, its traded
value and its number of trades. A \emph{time bar}\index{time bar} covers a fixed interval of wall-clock time; a \emph{tick
bar}\index{tick bar} a fixed number of trades; a \emph{volume bar}\index{volume bar} closes on the trade at which its cumulative
volume reaches a threshold, and a \emph{dollar bar}\index{dollar bar} on the trade at which its cumulative traded value does.
The \emph{volume-weighted average price}\index{volume-weighted average price} (VWAP) of a set of trades is $\sum_i p_iq_i/\sum_i q_i$.
\end{definition}

\begin{definition}[Imbalance bar]\label{def:rs:market-data-for-research:imbalance}
An \emph{imbalance bar}\index{imbalance bar} closes when the absolute sum of the trade signs since its start, $|\theta| =
|\sum b_i|$, reaches a threshold set from the \omterm{def:rs:market-data-for-research:bar}{bars} already closed: the expected number of trades per \omterm{def:rs:market-data-for-research:bar}{bar} times the expected
absolute imbalance per trade (López de Prado's tick-imbalance \omterm{def:rs:market-data-for-research:bar}{bar}), floored at the square root of the expected number of trades.
\end{definition}

\begin{omfigure}
\begin{tikzpicture}[font=\footnotesize,x=1.25cm,y=0.22cm]
\draw[->,>=Stealth] (0,0) -- (10.4,0) node[right] {time};
\foreach \x/\h in {0.3/2,0.5/1,0.7/3,0.9/2,1.1/1,1.3/2,2.6/1,4.2/1,5.8/2,6.1/3,6.3/2,6.5/1,6.7/2,6.9/3,7.1/1,8.8/1,9.5/2}
  \draw[omDef,line width=1.6pt] (\x,0) -- (\x,\h*2);
\foreach \x in {2,4,6,8} \draw[black!60,dashed] (\x,-3.5) -- (\x,8.5);
\draw[omThm,very thick] (0.7,-8.5) -- (0.7,-4.5) (2.6,-8.5) -- (2.6,-4.5) (6.1,-8.5) -- (6.1,-4.5) (6.9,-8.5) -- (6.9,-4.5);
\node[black!70,anchor=west] at (0,10) {time clock: a bar every two units, whatever happens};
\node[omThm,anchor=west] at (0,-10.5) {volume clock: a bar each time six units of volume have traded};
\node[omDef,anchor=west] at (6.6,10) {trades (height = size)};
\end{tikzpicture}

\omcaption{The same seventeen trades cut by two clocks. The time clock closes five equal intervals (dashed): three nearly
empty, one holding the early burst and one the burst between $t = 6$ and 7. The volume clock (solid, below) closes four
\omterm{def:rs:market-data-for-research:bar}{bars}: two short ones inside the bursts, ending at 0.7 and 6.9, and two long ones that absorb the quiet stretches; the last
three trades wait for the next \omterm{def:rs:market-data-for-research:bar}{bar}.}\label{fig:rs:market-data-for-research:schematic}
\end{omfigure}

The floor is needed in practice. When buys and sells balance, the expected imbalance per trade is near zero and so is the
threshold, and every trade would close a \omterm{def:rs:market-data-for-research:bar}{bar}; the square root of $\E[T]$ is the level a random walk of signs first reaches after
about $\E[T]$ trades.

Why sample by activity rather than by time? Because the variance of a price change grows with the amount of trading behind
it, and trading is uneven.

\begin{proposition}[Time bars are a variance mixture]\label{prop:rs:market-data-for-research:mixture}
Suppose that, given the volume $V$ traded in a \omterm{def:rs:market-data-for-research:bar}{bar}, the \omterm{def:rs:market-data-for-research:bar}{bar}'s log return is $\mathcal N(0, \sigma^2V)$. Then the kurtosis of
\omterm{def:rs:market-data-for-research:bar}{time-bar} returns is
\[
\frac{\E[r^4]}{\E[r^2]^2} = \frac{3\,\E[V^2]}{\E[V]^2} = 3\bigl(1 + \mathrm{CV}(V)^2\bigr),
\]
with $\mathrm{CV}$ the coefficient of variation of the \omterm{def:rs:market-data-for-research:bar}{bars}' volumes, while \omterm{def:rs:market-data-for-research:bar}{bars} of equal volume have Gaussian returns.
\end{proposition}

\begin{proof}
$\E[r^2] = \sigma^2\E[V]$ and $\E[r^4] = 3\sigma^4\E[V^2]$ by conditioning on $V$. A \omterm{def:rs:market-data-for-research:bar}{volume bar} has $V$ constant, up to the last
trade's overshoot.
\end{proof}

Clark (1973) proposed this subordination of prices to trading volume to explain heavy-tailed daily returns, and Ané and Geman
(2000) found that returns sampled on a clock of trade counts were close to Gaussian. On the simulated day the one-minute \omterm{def:rs:market-data-for-research:bar}{bars}'
volumes have a coefficient of variation of 0.94, so the proposition predicts a kurtosis of 5.67; the one-minute midquote returns
have 5.44, and the close-to-close returns 5.47. The \omterm{def:rs:market-data-for-research:bar}{bars} cut by the other clocks, with the same number of \omterm{def:rs:market-data-for-research:bar}{bars} in the day (about
390), have kurtoses of 3.58 (tick), 4.04 (volume) and 3.92 (dollar): not 3, because in the simulator volatility follows the
activity level and not volume exactly, but much closer. The event clocks spend their \omterm{def:rs:market-data-for-research:bar}{bars} where the trading is
(\cref{fig:rs:market-data-for-research:clocks}): between 2 and 42 \omterm{def:rs:market-data-for-research:bar}{bars} per quarter-hour, against a steady 15 for the time clock.

\begin{omfigure}
\begin{tikzpicture}
\begin{axis}[width=11cm,height=5cm,xlabel={minutes after the open},ylabel={bars per quarter-hour},
  tick label style={font=\small},label style={font=\small},xmin=0,xmax=390,ymin=0,ymax=45,grid=major,grid style={gray!25},
  legend columns=3,legend style={font=\footnotesize,at={(0.5,-0.32)},anchor=north,draw=none,fill=none,
  /tikz/every even column/.append style={column sep=8pt}},table/col sep=comma]
\addplot[black!60,very thick,const plot mark mid] table[x=minute,y=time]{figdata/research/02-market-data-for-research/clocks.csv};
\addplot[omDef,very thick,const plot mark mid] table[x=minute,y=volume]{figdata/research/02-market-data-for-research/clocks.csv};
\addplot[omThm,thick,dashed,const plot mark mid] table[x=minute,y=tick]{figdata/research/02-market-data-for-research/clocks.csv};
\addplot[black,thin,sharp plot] coordinates {(210,0) (210,45)};
\legend{time bars,volume bars,tick bars}
\end{axis}
\end{tikzpicture}

\omcaption{\omterm{def:rs:market-data-for-research:bar}{Bars} closed per quarter-hour on the simulated day, each clock cutting about 390 \omterm{def:rs:market-data-for-research:bar}{bars}. The time clock closes 15 \omterm{def:rs:market-data-for-research:bar}{bars}
in every quarter-hour; the volume and tick clocks follow the activity: busy after the open, quiet in the middle of the day,
busy around the news burst at minute 210 (vertical line) and before the close. \omterm{def:rs:market-data-for-research:bar}{Dollar bars} coincide with \omterm{def:rs:market-data-for-research:bar}{volume bars} at this
scale. Data: \texttt{firm.tape}, the chapter's day, seeded.}\label{fig:rs:market-data-for-research:clocks}
\end{omfigure}

\section{What each representation destroys}

\begin{definition}[Midquote series]\label{def:rs:market-data-for-research:midquote}
A \emph{midquote series}\index{midquote series} is the mid price $\tfrac12(b_t + a_t)$ sampled at chosen times, each value the
last mid at or before the sampling time (an as-of sample).
\end{definition}

A trade price is the mid plus or minus half the spread, depending on who initiated the trade. Sampled too finely, it bounces.

\begin{proposition}[The bounce adds half a squared spread per return]\label{prop:rs:market-data-for-research:bounce}
Let the trade price be $p = m + \tfrac s2 b$, with a constant spread $s$ and trade signs $b = \pm1$ independent of each other and
of the mid. Between two trades, $\E[(\Delta p)^2] = \E[(\Delta m)^2] + s^2/2$.
\end{proposition}

\begin{proof}
$\Delta p = \Delta m + \tfrac s2(b_2 - b_1)$, and $\E[(b_2 - b_1)^2] = 2$ when the signs are independent with mean zero.
\end{proof}

This is Roll's model, whose estimator of the spread Book 4, chapter 21, derives; the same chapter's signature plot draws the
realised variance against the sampling interval. Here the plot answers the hook
(\cref{fig:rs:market-data-for-research:signature}). At one second the trade-price variance of the day is 0.83 (in squared
per cent) and the midquote's 0.23; at five minutes they are 0.84 and 0.85. The trade price overstates at high frequency, as the
proposition says. The midquote understates, for a reason the proposition does not see: in the simulator, as in real markets,
the mid catches up with the efficient price in steps, a stale quote at a time, so its short-horizon returns are positively
autocorrelated and their squares add up to less than the variance over longer horizons. Neither series is the efficient price.

\begin{omfigure}
\begin{tikzpicture}
\begin{semilogxaxis}[width=10cm,height=5.2cm,xlabel={sampling interval (seconds)},ylabel={realised variance (\%$^2$)},
  tick label style={font=\small},label style={font=\small},xmin=1,xmax=900,ymin=0,ymax=1.2,grid=major,grid style={gray!25},
  legend columns=2,legend style={font=\footnotesize,at={(0.5,-0.32)},anchor=north,draw=none,fill=none,
  /tikz/every even column/.append style={column sep=8pt}},table/col sep=comma]
\addplot[omThm,very thick,mark=*,mark size=1.2pt] table[x=seconds,y=trade]{figdata/research/02-market-data-for-research/signature.csv};
\addplot[omDef,very thick,mark=square*,mark size=1.2pt] table[x=seconds,y=mid]{figdata/research/02-market-data-for-research/signature.csv};
\legend{last trade price,midquote}
\end{semilogxaxis}
\end{tikzpicture}

\omcaption{Signature plot of the simulated day: realised variance of log returns against the sampling interval, from trade
prices and from midquotes. At one second, 0.83 against 0.23; from about a minute on the two agree, and the noise of a
single day's estimate dominates. Data: \texttt{firm.tape}, the chapter's day, seeded.}\label{fig:rs:market-data-for-research:signature}
\end{omfigure}

The bounce is one loss among several, and a research store should know which each representation suffers:
\begin{itemize}
\item \emph{Trade prices} carry the bounce, and trades flagged with special conditions (odd lots, out-of-sequence reports,
auction prints, late corrections: Book 1, chapter 28) mixed with regular ones. An opening auction print is a price at which a
large volume traded at one instant; treating it as one more trade distorts the first \omterm{def:rs:market-data-for-research:bar}{bar} of every day.
\item \emph{Midquotes} cannot see what traded or how much, and at a wide spread the mid can move without any trade.
\item \emph{\omterm{def:rs:market-data-for-research:bar}{Bars}} lose the order of events inside the \omterm{def:rs:market-data-for-research:bar}{bar}, the book, and every event between samples; the high and the low keep a
trace of the path, which chapter~7's range estimators exploit.
\item \emph{\omterm{def:rs:market-data-for-research:bar}{Time bars} on several instruments} are aligned in time but not in information: an illiquid instrument's last trade
may be minutes old (non-synchronous trading and the Epps effect, Book 4, chapter 21).
\item \emph{Any series stamped with one clock} hides the difference between the exchange's timestamp and the moment the data
arrived; a backtest must use the second (chapter~3).
\end{itemize}

\section{Continuous futures series}

A futures contract expires, and its successor trades at a different price. A history long enough for research must splice
contracts, and the splice is a choice.

\begin{definition}[Continuous futures series, back-adjustment, ratio adjustment]\label{def:rs:market-data-for-research:continuous}
A \emph{continuous futures series}\index{continuous futures series} follows the contract a position would hold under a roll rule:
the nearby contract until a roll date some days before its last trading day, then the next contract. \emph{Back-adjustment}\index{back-adjustment}
adds, to every price before a roll, the gap between the new and the old contract on the roll date, so that the series has no
jump at the roll and its latest prices are real. \emph{Ratio adjustment}\index{ratio adjustment} multiplies every earlier price by
the ratio of the new contract's price to the old one's on the roll date instead.
\end{definition}

\begin{proposition}[What each adjustment preserves]\label{prop:rs:market-data-for-research:adjust}
\omterm{def:rs:market-data-for-research:continuous}{Ratio adjustment} preserves the percentage returns of the rolled position on every day; \omterm{def:rs:market-data-for-research:continuous}{back-adjustment} preserves its price
changes, the profit of one contract. Neither preserves the price levels of the past, and back-adjusted prices can be negative.
\end{proposition}

\begin{proof}
Between two roll dates both adjustments act on the held contract's prices by a constant: an additive constant leaves differences
unchanged, a multiplicative one leaves ratios unchanged. Across a roll date the adjustment is chosen so that the new contract's
price change (or return) is recorded, not the jump between contracts. Additive constants accumulate the roll gaps and can exceed
a past price.
\end{proof}

The crude oil futures settlements published by the US Energy Information Administration make the choice concrete
(\cref{fig:rs:market-data-for-research:wti}). From 2 January 2015 to 5 April 2024 the nearby contract went from USD~52.69 to
USD~86.91 a barrel, up 65\%. A long position rolled five trading days before each of the 111 expiries of the period, the way a
fund holding the nearby contract must, lost 27\%: most of those years were in contango, and each roll sold a cheaper contract to
buy a dearer one (the roll yield of One Quant Book 3, chapter 10). The ratio-adjusted series tells the fund's story; its first price is
USD~118.33. The back-adjusted series starts at USD~70.23 and, because the backwardation of 2021--2023 made the later gaps
negative, reaches $-5.89$ on 21 April 2020, a day when the contract actually held (June 2020) settled at USD~11.57. The nearby
series itself shows $-37.63$ on 20 April 2020, the settlement of an expiring contract no rolled position held.

\begin{omfigure}
\begin{tikzpicture}
\begin{axis}[width=11cm,height=5.6cm,xlabel={year},ylabel={USD per barrel},tick label style={font=\small},
  label style={font=\small},xmin=2015,xmax=2024.3,ymin=-40,ymax=130,grid=major,grid style={gray!25},
  xtick={2015,2016,2017,2018,2019,2020,2021,2022,2023,2024},xticklabels={2015,2016,2017,2018,2019,2020,2021,2022,2023,2024},
  scaled x ticks=false,legend columns=3,legend style={font=\footnotesize,at={(0.5,-0.3)},anchor=north,draw=none,fill=none,
  /tikz/every even column/.append style={column sep=8pt}},table/col sep=comma]
\addplot[black!45,thin] table[x=year,y=front]{figdata/research/02-market-data-for-research/wti.csv};
\addplot[omDef,thick] table[x=year,y=ratio]{figdata/research/02-market-data-for-research/wti.csv};
\addplot[omThm,thick,dashed] table[x=year,y=back]{figdata/research/02-market-data-for-research/wti.csv};
\legend{nearby (as published),ratio-adjusted,back-adjusted}
\end{axis}
\end{tikzpicture}

\omcaption{WTI crude oil futures, 2015 to April 2024: the nearby contract as published (with $-37.63$ on 20 April 2020), and
continuous series rolling five trading days before each expiry, ratio-adjusted (starting at 118.33) and back-adjusted (starting
at 70.23, negative in April 2020). All three end at the same price. Data: US Energy Information Administration, NYMEX settlements,
contracts 1 and 2.}\label{fig:rs:market-data-for-research:wti}
\end{omfigure}

\begin{remark}[Which series for which question]\label{rem:rs:market-data-for-research:which}
Returns of a rolled strategy: ratio-adjusted. Profit of a fixed number of contracts: back-adjusted differences. A signal on
price levels (a moving average of the price, a breakout above a past high): neither without care, since both rewrite the past
at every roll; compute the signal on the contract held at the time. The shape of the curve (carry, Book 8): the raw contracts,
never an adjusted series.
\end{remark}

\section{Volumes, storage and formats}

One simulated instrument-day of 1.56 million messages at 36 bytes (the size of an ITCH add-order message) is 56~MB before
compression; the dated box shows what a whole exchange-day weighs. Three rules keep a research store usable.

\begin{method}[Storing market data for research]\label{met:rs:market-data-for-research:store}
\begin{enumerate}
\item Keep the raw messages as received, immutable, with both timestamps, partitioned by date and instrument; every derived
table (books, \omterm{def:rs:market-data-for-research:bar}{bars}, features) is rebuilt from them by versioned code.
\item Store derived tables in a columnar, compressed format and partition them the way they are read: by date for
cross-sectional work, by instrument for time-series work.
\item Record for every derived table the code version, the raw partitions it read and its parameters (the clock, the
threshold, the roll rule), so that a result can name the exact data it used (chapter~29).
\end{enumerate}
\end{method}

\section{Tutorial: one day, four clocks}\label{tut:rs:market-data-for-research:clocks}

\begin{tutorial}
\textbf{Goal.} Simulate a trading day, cut it by four clocks and measure what each clock and each price does to the returns.
\textbf{End state:} \cref{fig:rs:market-data-for-research:clocks,fig:rs:market-data-for-research:signature}; kurtoses 5.47 (time),
3.58 (tick), 4.04 (volume), 3.92 (dollar); the one-second variance ratio of 3.6.
\begin{tutsteps}
\item \textbf{The day.} \texttt{simulate(TapeConfig(seconds=23\,400, u\_shape=1.5, news\_at=12\,600, seed=3))} returns the
messages, the trades, the top of book after every message and the efficient price (about 10 seconds).
\item \textbf{Threshold \omterm{def:rs:market-data-for-research:bar}{bars}.} A \omterm{def:rs:market-data-for-research:bar}{bar} closes on the trade that crosses the threshold; trades are never split between \omterm{def:rs:market-data-for-research:bar}{bars}.
\omcode{code/firm/bars/firm_bars.py}{65}{90}{Volume and dollar bars close on the crossing trade.}\label{lst:rs:market-data-for-research:cuts}
\item \textbf{Two prices, one day.} Sample the last trade and the midquote on grids of 1 to 900 seconds and sum the squared
log returns; compare the kurtosis with the mixture prediction from the \omterm{def:rs:market-data-for-research:bar}{bars}' volumes.
\omcode{code/research/02-market-data-for-research/python/rs_marketdata.py}{45}{78}{Realised variance from trades and from mids; the mixture kurtosis.}\label{lst:rs:market-data-for-research:rv}
\item \textbf{Run} \texttt{summary()} and \texttt{fig\_marketdata.py}.
\end{tutsteps}
\textbf{What to change next.} Set \texttt{act\_vol=0} (no random activity) and watch the \omterm{def:rs:market-data-for-research:bar}{time bars}' kurtosis fall towards the
mixture value of the U-shape alone; cut \omterm{def:rs:market-data-for-research:imbalance}{imbalance bars} and follow their length through the day (exercise~7).
\end{tutorial}

\section{Build: the market simulator and the bar builders}\label{bld:rs:market-data-for-research:tape}

\begin{build}
\bfield{Purpose} The book's standard intraday data, with the truth attached: \texttt{firm.tape} writes an order-by-order feed
from a known mechanism (efficient price, informed and noise flow, stale-quote cancellations); \texttt{firm.bars} turns trades
into \omterm{def:rs:market-data-for-research:bar}{bars}. Chapters 8--10 measure predictors on the tape, chapter~18 replays it, chapter~19 trades against it.
\bfield{Interface} \texttt{TapeConfig} (duration, tick, levels, rates, Hawkes and metaorder parameters, efficient-price jump
rate, news window, intraday profile, activity process, seed); \texttt{simulate(cfg, v\_path)} returning a \texttt{Tape} of
structured arrays \texttt{msgs}, \texttt{trades} (with true signs and an informed flag), \texttt{top}, and the paths \texttt{v\_t},
\texttt{v}, \texttt{act}; \texttt{simulate\_pair(cfg, latency)}; \texttt{Book().apply(msg)}. \texttt{time\_bars},
\texttt{tick\_bars}, \texttt{volume\_bars}, \texttt{dollar\_bars}, \texttt{imbalance\_bars}, \texttt{vwap}, \texttt{asof},
\texttt{continuous(c1, c2, expiry, days\_before)}.
\bfield{Rules} Deterministic for a seed. Rates piecewise constant between events except the Hawkes intensity, simulated exactly
by thinning. The opening book is a snapshot at time zero. Prices are integer ticks, sizes integer shares.
\bfield{Acceptance tests} \texttt{code/firm/tape/tests/} and \texttt{code/firm/bars/tests/}: identical output for a seed;
the messages rebuild the recorded top of book at every step and the book never crosses; executions equal trades; the mid tracks
the efficient price; pooled over six seeds, imbalance predicts the next mid move and informed trades lose money for the resting
side while noise trades pay it; \omterm{def:rs:market-data-for-research:bar}{bars} close on the crossing trade; a hand-checked roll.
\bfield{Stretch} Several instruments sharing a factor in their efficient prices; auctions at the open and close; a C++20 port
of the event loop.
\end{build}

\begin{omsources}
\item P.~K.~Clark, ``A subordinated stochastic process model with finite variance for speculative prices'',
\emph{Econometrica} 41(1), 1973.
\item T.~Ané and H.~Geman, ``Order flow, transaction clock, and normality of asset returns'', \emph{Journal of Finance} 55(5),
2000.
\item D.~Easley, M.~López de Prado and M.~O'Hara, ``The volume clock: insights into the high-frequency paradigm'',
\emph{Journal of Portfolio Management} 39(1), 2012.
\item M.~López de Prado, \emph{Advances in Financial Machine Learning}, Wiley, 2018, chapter 2 (information-driven bars).
\item R.~Roll, ``A simple implicit measure of the effective bid--ask spread in an efficient market'', \emph{Journal of Finance}
39(4), 1984.
\item Nasdaq, \emph{TotalView-ITCH 5.0} specification, and the sample files on \texttt{emi.nasdaq.com}.
\item US Energy Information Administration, NYMEX futures prices, crude oil contracts 1--4 (daily).
\end{omsources}

\section{Exercises}

\begin{exercise}[$\star$]\label{exo:rs:market-data-for-research:1}
Trades of 300 shares at 20.00, 500 at 20.02 and 200 at 19.99. What is their VWAP?
\end{exercise}

\begin{exercise}[$\star$]\label{exo:rs:market-data-for-research:2}
A stock at USD~20 has a one-cent spread and trades 20\,000 times a day with a true daily volatility of 2\%. By
\cref{prop:rs:market-data-for-research:bounce}, what daily volatility does the sum of squared trade-to-trade log returns report?
\end{exercise}

\begin{exercise}[$\star$]\label{exo:rs:market-data-for-research:3}
A feed averages 40 bytes a message and an active instrument 1.5 million messages a day. How much raw data do 5\,000 such
instruments produce in a year of 252 days?
\end{exercise}

\begin{exercise}[$\star\star$]\label{exo:rs:market-data-for-research:4}
One-minute volumes have a coefficient of variation of 0.8. What kurtosis does \cref{prop:rs:market-data-for-research:mixture}
predict for one-minute returns? What coefficient of variation gives a kurtosis of 6?
\end{exercise}

\begin{exercise}[$\star\star$]\label{exo:rs:market-data-for-research:5}
A nearby contract settles at 70, 71, 72 on three days and the next contract at 73, 74, 76; the position rolls at the close of
the second day. Give the held series, the back-adjusted and the ratio-adjusted series.
\end{exercise}

\begin{exercise}[$\star\star$]\label{exo:rs:market-data-for-research:6}
A stock traded at USD~10 in 2015 and at USD~100 in 2025 with the same traded value every day. A \omterm{def:rs:market-data-for-research:bar}{volume-bar} threshold fixed in
shares gives how many \omterm{def:rs:market-data-for-research:bar}{bars} a day in 2015 relative to 2025? And a \omterm{def:rs:market-data-for-research:bar}{dollar-bar} threshold?
\end{exercise}

\begin{exercise}[$\star\star\star$]\label{exo:rs:market-data-for-research:7}
\emph{Coding.} Cut the chapter's day into \omterm{def:rs:market-data-for-research:imbalance}{imbalance bars}, starting from an expected 142 trades a \omterm{def:rs:market-data-for-research:bar}{bar}. How many \omterm{def:rs:market-data-for-research:bar}{bars} close, how
many trades do they hold on average, and what is the kurtosis of their midquote returns? Why do the \omterm{def:rs:market-data-for-research:bar}{bars} shrink through the day?
\end{exercise}

\begin{exercise}[$\star\star\star$]\label{exo:rs:market-data-for-research:8}
\emph{Find the flaw.} ``Our carry signal is the annualised log ratio of the back-adjusted second-month series to the
back-adjusted front-month series, and our trend signal the percentage change of the back-adjusted front-month series over
twelve months.''
\end{exercise}

\section{Problem: Which Clock?}

\begin{problem}[{Weekend problem --- one day, four clocks and two prices}]\label{pb:rs:market-data-for-research:1}
The chapter's simulated day: 6.5 hours, a U-shaped profile, a random activity level and a news burst at 12\,600 seconds, seed 3.

\medskip\noindent\textbf{Part I --- The feed.}
\begin{enumerate}
\item How many messages and trades does the day have, and how many messages per trade?
\item What share of the messages are additions, cancellations and executions, roughly, and why are additions and cancellations
nearly equal?
\item What is the raw size of the day at 36 bytes a message?
\item Why does a research store keep the messages rather than the \omterm{def:rs:market-data-for-research:bar}{bars}?
\end{enumerate}

\medskip\noindent\textbf{Part II --- Clocks.}
\begin{enumerate}[resume]
\item How many \omterm{def:rs:market-data-for-research:bar}{bars} does each clock cut, with thresholds chosen for about 390?
\item What are the kurtoses of close-to-close returns for the four clocks?
\item What is the coefficient of variation of the one-minute \omterm{def:rs:market-data-for-research:bar}{bars}' volumes, and what kurtosis does the mixture proposition
predict for them?
\item Why are the event clocks' kurtoses above 3?
\item Where in the day do \omterm{def:rs:market-data-for-research:bar}{volume bars} concentrate, and why?
\end{enumerate}

\medskip\noindent\textbf{Part III --- Prices.}
\begin{enumerate}[resume]
\item What are the realised variances from trades and from mids at 1, 5, 60 and 300 seconds?
\item Which effect makes the trade series too high at one second?
\item Which makes the mid series too low?
\item From which interval on do the two agree within 5\%?
\item What would you use to estimate the day's variance from the one-second data (Book 4, chapter 21)?
\end{enumerate}

\medskip\noindent\textbf{Part IV --- Judgement.}
\begin{enumerate}[resume]
\item Which clock would you use for a daily volatility estimate, and which for a model of short-term price moves?
\item What does a one-minute \omterm{def:rs:market-data-for-research:bar}{time bar} lose that a \omterm{def:rs:market-data-for-research:bar}{volume bar} keeps?
\item What does a \omterm{def:rs:market-data-for-research:bar}{volume bar} lose that a \omterm{def:rs:market-data-for-research:bar}{time bar} keeps?
\item State the \emph{named result}: the kurtosis of one-minute \omterm{def:rs:market-data-for-research:bar}{time-bar} returns against \omterm{def:rs:market-data-for-research:bar}{dollar-bar} returns on the day, and the
trade-to-mid variance ratio at one second.
\item What in the simulator makes the \omterm{def:rs:market-data-for-research:bar}{time bars} heavy-tailed, and what would remove it?
\item In one sentence: which representation is ``the'' price?
\end{enumerate}
\end{problem}

\section{Interview questions}

\begin{interviewq}[$\star$ \iqroles{researcher, trader}]\label{iq:rs:market-data-for-research:1}
Why might you sample prices by volume rather than by time?
\end{interviewq}

\begin{interviewq}[$\star\star$ \iqroles{researcher}]\label{iq:rs:market-data-for-research:2}
Your realised volatility from one-second trade prices is double the one from five-minute prices. What is going on?
\end{interviewq}

\begin{interviewq}[$\star\star$ \iqroles{researcher, trader}]\label{iq:rs:market-data-for-research:3}
How would you build a \omterm{def:rs:market-data-for-research:continuous}{continuous futures series} for a trend-following backtest, and what goes wrong with back-adjusted prices?
\end{interviewq}

\begin{interviewq}[$\star\star$ \iqroles{developer, researcher}]\label{iq:rs:market-data-for-research:4}
Design the storage of a year of order-by-order data for 5\,000 instruments for research use.
\end{interviewq}

\begin{interviewq}[$\star\star$ \iqroles{researcher, mle}]\label{iq:rs:market-data-for-research:5}
A colleague's model uses one-minute \omterm{def:rs:market-data-for-research:bar}{bars} of fifty stocks. What can go wrong because the \omterm{def:rs:market-data-for-research:bar}{bars} are \omterm{def:rs:market-data-for-research:bar}{time bars}?
\end{interviewq}

\begin{interviewq}[$\star\star\star$ \iqroles{researcher}]\label{iq:rs:market-data-for-research:6}
Show that if returns are normal given the volume traded, \omterm{def:rs:market-data-for-research:bar}{time-bar} returns have kurtosis $3(1 + \mathrm{CV}^2)$.
\end{interviewq}
